\documentclass[../main.tex]{subfiles}
\ifSubfilesClassLoaded{\setcounter{chapter}{10}}{}
\begin{document}

\chapter{Struktur Data Bentukan: \texttt{struct}, \texttt{union}, \texttt{enum}, dan \texttt{typedef}}

\begin{subcpmk}
  \item Sub-CPMK 5.3: \texttt{struct} --- definisi, akses anggota, nested struct, array of struct, pointer ke struct (\texttt{->}), \texttt{typedef}
  \item Sub-CPMK 5.4: \texttt{union}, \texttt{enum}, bit fields (pengenalan)
\end{subcpmk}

\SesiRPS{11}{5.3--5.4}{$\sim$170 menit}{CPMK-5}

% ============================================================
\section{Mengapa Tipe Data Bentukan?}
% ============================================================

Tipe data dasar (\texttt{int}, \texttt{double}, \texttt{char}) menyimpan satu nilai. Namun data di dunia nyata sering kali terdiri dari beberapa informasi terkait --- misalnya data mahasiswa memiliki NIM (string), nama (string), dan IPK (double). \texttt{struct} memungkinkan kita mengelompokkan data-data terkait menjadi satu \textbf{tipe data baru}, sehingga kode lebih terorganisir, terbaca, dan mudah dikelola dibandingkan menggunakan banyak variabel terpisah.

Selain \texttt{struct}, C juga menyediakan \texttt{union} (menyimpan satu dari beberapa tipe di lokasi memori yang sama), \texttt{enum} (mendefinisikan sekumpulan konstanta bernama), dan \texttt{typedef} (membuat alias untuk tipe data). Konsep-konsep ini merupakan fondasi untuk membangun abstraksi data yang lebih kompleks.

% ============================================================
\section{Mendefinisikan dan Menggunakan \texttt{struct}}
% ============================================================

\begin{lstlisting}[language=C,caption=Definisi dan penggunaan \texttt{struct}]
#include <stdio.h>
#include <string.h>

// Definisi struct
struct Mahasiswa {
    char nim[16];
    char nama[50];
    double ipk;
    int semester;
};

int main(void) {
    // Deklarasi dan inisialisasi
    struct Mahasiswa mhs1 = {"12345678", "Budi Santoso", 3.75, 4};
    
    // Akses anggota dengan operator titik (.)
    printf("NIM    : %s\n", mhs1.nim);
    printf("Nama   : %s\n", mhs1.nama);
    printf("IPK    : %.2f\n", mhs1.ipk);
    printf("Semester: %d\n", mhs1.semester);
    
    // Mengubah nilai anggota
    mhs1.ipk = 3.80;
    mhs1.semester = 5;
    
    return 0;
}
\end{lstlisting}

Berikut adalah contoh penggunaan \texttt{struct} yang langsung memproses data pegawai (nama, golongan, jumlah anak, upah) --- konversi dari Pascal \texttt{record\_upah.pas}.

\lstinputlisting[language=C,caption={\texttt{contoh\_code/c/record\_upah.c} --- \texttt{struct} pegawai langsung tanpa fungsi (konversi Pascal)}]{../contoh_code/c/record_upah.c}

% ============================================================
\section{\texttt{typedef} untuk Menyederhanakan Nama Tipe}
% ============================================================

Menulis \texttt{struct Mahasiswa} berulang kali cukup panjang. \texttt{typedef} membuat alias sehingga kita bisa menulis nama tipe yang lebih singkat. Ini adalah konvensi yang sangat umum di C --- hampir semua proyek C profesional menggunakan \texttt{typedef} untuk struct.

\begin{lstlisting}[language=C,caption=\texttt{typedef} untuk struct]
#include <stdio.h>

typedef struct {
    char nim[16];
    char nama[50];
    double ipk;
} Mahasiswa;  // sekarang bisa langsung tulis Mahasiswa

int main(void) {
    Mahasiswa m = {"12345678", "Siti Aminah", 3.45};
    printf("%s - %s - IPK: %.2f\n", m.nim, m.nama, m.ipk);
    return 0;
}
\end{lstlisting}

% ============================================================
\section{Array of Struct}
% ============================================================

Menggabungkan array dan struct memungkinkan kita menyimpan kumpulan data terstruktur, seperti daftar mahasiswa, inventory barang, atau database sederhana.

\begin{lstlisting}[language=C,caption=Array of struct --- daftar mahasiswa]
#include <stdio.h>
#include <string.h>

typedef struct {
    char nim[16];
    char nama[50];
    double ipk;
} Mahasiswa;

void cetak_daftar(const Mahasiswa daftar[], int n) {
    printf("%-12s %-20s %s\n", "NIM", "Nama", "IPK");
    printf("--------------------------------------\n");
    for (int i = 0; i < n; i++) {
        printf("%-12s %-20s %.2f\n",
               daftar[i].nim, daftar[i].nama, daftar[i].ipk);
    }
}

int main(void) {
    Mahasiswa daftar[] = {
        {"11111111", "Andi", 3.50},
        {"22222222", "Budi", 3.75},
        {"33333333", "Cici", 3.90},
    };
    int n = sizeof(daftar) / sizeof(daftar[0]);
    
    cetak_daftar(daftar, n);
    return 0;
}
\end{lstlisting}

\lstinputlisting[language=C,caption={\texttt{contoh\_code/c/daf\_pegawai\_array.c} --- Array of struct pegawai dengan fungsi input/output (konversi Pascal)}]{../contoh_code/c/daf_pegawai_array.c}

% ============================================================
\section{Pointer ke Struct dan Operator \texttt{->}}
% ============================================================

Ketika kita memiliki pointer ke struct, kita menggunakan operator \textbf{arrow} (\texttt{->}) untuk mengakses anggotanya. \texttt{p->nama} setara dengan \texttt{(*p).nama} tetapi lebih terbaca. Pointer ke struct sangat penting untuk fungsi yang harus mengubah data struct dan untuk struktur data dinamis (linked list, di Bab selanjutnya).

\begin{lstlisting}[language=C,caption=Pointer ke struct dan operator \texttt{->}]
#include <stdio.h>

typedef struct {
    char nama[50];
    int umur;
} Orang;

void ulang_tahun(Orang *p) {
    p->umur++;  // setara dengan (*p).umur++
    printf("%s sekarang berumur %d tahun\n", p->nama, p->umur);
}

int main(void) {
    Orang andi = {"Andi", 20};
    
    printf("Sebelum: %s, %d tahun\n", andi.nama, andi.umur);
    ulang_tahun(&andi);  // kirim alamat struct
    printf("Sesudah: %s, %d tahun\n", andi.nama, andi.umur);
    
    return 0;
}
\end{lstlisting}

\lstinputlisting[language=C,caption={\texttt{contoh\_code/c/record\_prosedur.c} --- \texttt{struct} pegawai dengan fungsi \texttt{input} dan \texttt{cetak} (konversi Pascal)}]{../contoh_code/c/record_prosedur.c}

% ============================================================
\section{Nested Struct}
% ============================================================

\begin{lstlisting}[language=C,caption=Struct di dalam struct]
#include <stdio.h>

typedef struct {
    int hari, bulan, tahun;
} Tanggal;

typedef struct {
    char nama[50];
    Tanggal tgl_lahir;  // struct di dalam struct
} Pegawai;

int main(void) {
    Pegawai p = {"Dewi", {15, 8, 1995}};
    
    printf("Nama: %s\n", p.nama);
    printf("Lahir: %02d/%02d/%d\n",
           p.tgl_lahir.hari, p.tgl_lahir.bulan, p.tgl_lahir.tahun);
    
    return 0;
}
\end{lstlisting}

% ============================================================
\section{\texttt{union}}
% ============================================================

\texttt{union} mirip \texttt{struct}, tetapi semua anggotanya \textbf{berbagi lokasi memori yang sama}. Ukuran union adalah ukuran anggota terbesar. Hanya satu anggota yang ``aktif'' pada satu waktu --- menulis ke satu anggota menimpa yang lain. Union berguna untuk menghemat memori ketika hanya satu dari beberapa tipe data yang akan digunakan.

\begin{modultable}[]{Perbandingan \texttt{struct} dan \texttt{union}}
\begin{tabular}{|l|p{5cm}|p{5cm}|}
\hline
\textbf{Aspek} & \texttt{\textbf{struct}} & \texttt{\textbf{union}} \\
\hline
Memori & Alokasi untuk \textbf{semua} anggota & Alokasi untuk anggota \textbf{terbesar} saja \\
\hline
Akses & Semua anggota bisa diakses bersamaan & Hanya 1 anggota yang valid \\
\hline
Ukuran & Jumlah semua anggota + padding & Ukuran anggota terbesar \\
\hline
\end{tabular}
\end{modultable}

\begin{lstlisting}[language=C,caption=Penggunaan \texttt{union}]
#include <stdio.h>

typedef union {
    int bilangan;
    double desimal;
    char karakter;
} Nilai;

int main(void) {
    Nilai v;
    
    v.bilangan = 42;
    printf("int: %d\n", v.bilangan);
    
    v.desimal = 3.14;  // menimpa v.bilangan
    printf("double: %.2f\n", v.desimal);
    // printf("int sekarang: %d\n", v.bilangan);  // nilai tidak valid!
    
    printf("sizeof(Nilai) = %zu\n", sizeof(Nilai));  // = sizeof(double)
    
    return 0;
}
\end{lstlisting}

% ============================================================
\section{\texttt{enum} (Enumerasi)}
% ============================================================

\texttt{enum} mendefinisikan sekumpulan \textbf{konstanta bernama} (enumerator) yang secara default bernilai integer berurutan mulai dari 0. Enum meningkatkan keterbacaan kode karena mengganti ``angka ajaib'' (\textit{magic numbers}) dengan nama yang bermakna.

\begin{lstlisting}[language=C,caption=Penggunaan \texttt{enum}]
#include <stdio.h>

typedef enum {
    SENIN = 1,     // mulai dari 1 (bukan default 0)
    SELASA,        // otomatis 2
    RABU,          // 3
    KAMIS,         // 4
    JUMAT,         // 5
    SABTU,         // 6
    MINGGU         // 7
} Hari;

const char *nama_hari(Hari h) {
    const char *nama[] = {"", "Senin", "Selasa", "Rabu",
                          "Kamis", "Jumat", "Sabtu", "Minggu"};
    return nama[h];
}

int main(void) {
    Hari hari_ini = RABU;
    printf("Hari ini: %s (nilai enum: %d)\n",
           nama_hari(hari_ini), hari_ini);
    return 0;
}
\end{lstlisting}

% ============================================================
\section{Referensi sesi}
% ============================================================
\begin{itemize}[leftmargin=*]
  \item cppreference, \texttt{struct}. \url{https://en.cppreference.com/w/c/language/struct}
  \item Petani Kode, struct. \url{https://petanikode.com/tutorial/c}
  \item Programiz, \textit{C struct}. \url{https://www.programiz.com/c-programming/c-structures}
  \item Beej's Guide, \textit{Structs}. \url{https://beej.us/guide/bgc/html/index.html}
\end{itemize}

% ============================================================
\begin{aktivitas}
  \item Buat struct \texttt{Buku} (judul, penulis, tahun, harga) dan cetak datanya.
  \item Buat array of struct untuk menyimpan 5 buku, lalu cari buku dengan harga tertinggi.
  \item Buat fungsi yang menerima pointer ke struct dan mengubah datanya.
  \item \textbf{Latihan kerangka:} berikut konversi dari Pascal \texttt{prosedur\_upah\_latihan.pas}. Lengkapi fungsi \texttt{input} dan \texttt{cetak} yang masih kosong.
\end{aktivitas}

\lstinputlisting[language=C,caption={\texttt{contoh\_code/c/prosedur\_upah\_latihan.c} --- Latihan: fungsi upah dengan struct (kerangka konversi Pascal)}]{../contoh_code/c/prosedur_upah_latihan.c}

\begin{latihan}
  \item Apa perbedaan \texttt{m.nama} dan \texttt{p->nama}? Kapan masing-masing digunakan?
  \item Mengapa \texttt{union} menghemat memori dibanding \texttt{struct}?
  \item Kapan \texttt{enum} lebih baik daripada \texttt{\#define}?
  \item Apa fungsi \texttt{typedef} dan mengapa ia umum dipakai dengan \texttt{struct}?
\end{latihan}

\begin{asesmen}
\textbf{Tugas M9:} daftar statis \texttt{struct} (min.\ 5 rekaman) dengan pencarian berdasarkan NIM atau nama, dan tampilan terformat.
\end{asesmen}

\begin{checklist}
  \item Saya dapat mendefinisikan dan menggunakan \texttt{struct}
  \item Saya memahami operator \texttt{.} dan \texttt{->}
  \item Saya dapat menggunakan array of struct
  \item Saya memahami perbedaan \texttt{struct} dan \texttt{union}
  \item Saya dapat menggunakan \texttt{enum} dan \texttt{typedef}
\end{checklist}

\begin{rangkuman}
Bab ini membahas tipe data bentukan: \texttt{struct} (mengelompokkan data terkait), \texttt{typedef} (alias tipe), array of struct, pointer ke struct (operator \texttt{->}), nested struct, \texttt{union} (anggota berbagi memori), dan \texttt{enum} (konstanta bernama).

\textbf{Poin kunci:} \texttt{typedef} menyederhanakan nama tipe; \texttt{p->anggota} untuk pointer ke struct; \texttt{union} hanya satu anggota aktif; \texttt{enum} mengganti magic numbers.

\textbf{Kata kunci:} \texttt{struct}, \texttt{typedef}, \texttt{union}, \texttt{enum}, operator titik, operator arrow, nested struct
\end{rangkuman}

\end{document}
